\scanpage{357}
that \deleted{\unsure{already}} if $\unsure{\Pi|\pi}_0,\unsure{\Pi|\pi}_1$ are two discretifications \inserted{(\unsure{oriented})} of $\unsure{\Pi|\pi}$, which define the same prediscretification (\unsure{oriented}), then the same holds for \inserted{their \unsure{traces} on $\unsure{\Pi|\pi}'$,} $\unsure{\Pi|\pi}'_0$ \unsure{and} $\unsure{\Pi|\pi}'_1$, with respect to $\unsure{\Pi|\pi}'$. But already if $\unsure{\Pi|\pi}_0$ and $\unsure{\Pi|\pi}_1$ are $\unsure{\Pi|\pi}$-conjugate (i.e. define same discretification (\unsure{oriented}, \uline{\unsure{strict}})), this is not so clear --- \deleted{\unsure{unless}} but yes, for we will have $\unsure{\Pi|\pi}=\unsure{\Pi|\pi}'\cdot\unsure{\Pi|\pi}_0$ hence if $\unsure{\Pi|\pi}_1=\operatorname{int}(u)\unsure{\Pi|\pi}_0$ writing $u=u'u_0$ with $u'\in\unsure{\Pi|\pi}'$, $u_0\in\unsure{\Pi|\pi}_0$, we will have
\[
\unsure{\Pi|\pi}_1=\operatorname{int}(u')\operatorname{int}(u_0)\unsure{\Pi|\pi}_0=\operatorname{int}(u')\unsure{\Pi|\pi}_0,
\]
hence $\unsure{\Pi|\pi}'_1=\operatorname{int}(u')\unsure{\Pi|\pi}'_0$ hence $\unsure{\Pi|\pi}'_0$ and $\unsure{\Pi|\pi}'_1$ are conjugate.

Let us try to proceed in the same way in the general case of a $u\in\widehat{\unsure{\mathcal{C}|\mathcal{G}}(\unsure{\Pi|\pi}_0)}^+$, \inserted{such that $\unsure{\Pi|\pi}_1=u(\unsure{\Pi|\pi}_0)$,} writing if possible $u=u'u_0$, with $u'\in\widehat{\unsure{\mathcal{C}|\mathcal{G}}(\unsure{\Pi|\pi}_0,\unsure{\Pi|\pi}'_0)}^+$, and $u_0\in\unsure{\mathcal{C}|\mathcal{G}}(\unsure{\Pi|\pi}_0)^+$, where we define $\unsure{\mathcal{C}|\mathcal{G}}(\unsure{\Pi|\pi}_0,\unsure{\Pi|\pi}'_0)^+$ as the group of \unsure{i-automorphisms|l-automorphisms} \inserted{\unsure{direct|discrete}} of $\unsure{\Pi|\pi}_0$ which \uline{stabilize $\unsure{\Pi|\pi}'_0$}, and $\widehat{\unsure{\mathcal{C}|\mathcal{G}}}(\unsure{\Pi|\pi}_0,\unsure{\Pi|\pi}'_0)^+$ is its profinite \unsure{compactification}. We would then have $\unsure{\Pi|\pi}_1=u'(u_0(\unsure{\Pi|\pi}_0))=\dot{u}'(\unsure{\Pi|\pi}_0)$, where \deleted{\illegible{1}} $\dot{u}'$ denotes the image of $u'$ in $\widehat{\unsure{\mathcal{T}|\mathcal{C}}(\unsure{\Pi|\pi}'_0)}$. In any case, there is no problem if $\unsure{\Pi|\pi}'\deleted{\unsure{_0}}$ is a \deleted{\unsure{closed}} \uline{characteristic} subgroup of $\unsure{\Pi|\pi}$, for then we have a homomorphism

\scanpage{358}
\mspage{181}
discrete $\unsure{\mathcal{C}|\mathcal{G}}(\unsure{\Pi|\pi}_0)\longrightarrow\unsure{\mathcal{C}|\mathcal{G}}(\unsure{\Pi|\pi}'_0)$, defining \unsure{the|the} $\unsure{\widehat{\unsure{\mathcal{C}|\mathcal{G}}}(\widehat{\unsure{\Pi|\pi}}_0)}\longrightarrow\unsure{\widehat{\unsure{\mathcal{C}|\mathcal{G}}}(\widehat{\unsure{\Pi|\pi}}'_0)}$. Since the \unsure{open} characteristic subgroups of $\unsure{\Pi|\pi}$ are cofinal, we are reduced (for the questions of factorizability \inserted{of $D(\unsure{\Pi|\pi})\longrightarrow D(\unsure{\Pi|\pi}')$ into $P(\unsure{\Pi|\pi})\longrightarrow P(\unsure{\Pi|\pi}')$ and $A(\unsure{\Pi|\pi})\longrightarrow A(\unsure{\Pi|\pi}')$} and of bijectivity of the maps $\unsure{P}(\unsure{\Pi|\pi})\longrightarrow P(\unsure{\Pi|\pi}')$ and $A(\unsure{\Pi|\pi})\longrightarrow A(\unsure{\Pi|\pi}')$) we are reduced to the case where $\unsure{\Pi|\pi}'$ is a \uline{characteristic} subgroup. We have \unsure{indeed|seen} \inserted{then} \deleted{the} factorizability of $D(\unsure{\Pi|\pi})\longrightarrow D(\unsure{\Pi|\pi}')$ into $P(\unsure{\Pi|\pi})\longrightarrow P(\unsure{\Pi|\pi}')$, \inserted{\unsure{we want}} \unsure{to show|let us show} that this map is injective \unsure{...} \inserted{This} \deleted{\unsure{this means}}
\inserted{\textcircled{a}} means that the inverse image in $\widehat{\widehat{\unsure{\mathcal{C}|\mathcal{G}}}}(\widehat{\unsure{\Pi|\pi}}_0)$ of $\widehat{\unsure{\mathcal{C}|\mathcal{G}}}(\widehat{\unsure{\Pi|\pi}}'_0)^+$ is $\widehat{\unsure{\mathcal{C}|\mathcal{G}}}(\widehat{\unsure{\Pi|\pi}}_0)^+$, the surjectivity \inserted{\textcircled{b}} means that every element of $\widehat{\widehat{\unsure{\mathcal{C}|\mathcal{G}}}}(\widehat{\unsure{\Pi|\pi}}'_0)$ is congruent modulo $\widehat{\unsure{\mathcal{C}|\mathcal{G}}(\unsure{\Pi|\pi}'_0)}^+$ to an element of the image of $\widehat{\widehat{\unsure{\mathcal{C}|\mathcal{G}}}}(\unsure{\Pi|\pi}_0)$. The factorizability into $A(\unsure{\Pi|\pi})\longrightarrow A(\unsure{\Pi|\pi}')$ means that by \inserted{\textcircled{c}} the preceding map \deleted{\unsure{of}} $\widehat{\widehat{\unsure{\mathcal{C}|\mathcal{G}}}}(\unsure{\Pi|\pi}_0)\hookrightarrow\widehat{\widehat{\unsure{\mathcal{C}|\mathcal{G}}}}(\unsure{\Pi|\pi}'_0)$, (N.B. it is immediate that \unsure{it is} injective) \deleted{\unsure{is}} \inserted{sends}
\[
M(\unsure{\Pi|\pi}_0)=\operatorname{Norm}_{\unsure{\widehat{\widehat{\mathcal{C}}}|\widehat{\mathcal{C}}}(\unsure{\Pi|\pi}_0)}\bigl(\widehat{\unsure{\mathcal{C}|\mathcal{G}}}^{+}(\unsure{\Pi|\pi}_0)\bigr)
\]
into $M(\unsure{\Pi|\pi}'_0)$, the injectivity of $A(\unsure{\Pi|\pi})\longrightarrow A(\unsure{\Pi|\pi}')$ that

\scanpage{359}

\textcircled{d} we even have \unsure{that} $M(\unsure{\Pi|\pi})$ is
\unsure{exactly|exact\textsuperscript{ly}} the inverse image of $M(\unsure{\Pi|\pi}')$,
the \unsure{surjectivity} of
$A(\unsure{\Pi|\pi})\unsure{\twoheadrightarrow|\longrightarrow}A(\unsure{\Pi|\pi}')$ that
\inserted{\textcircled{e}} every element of
$\widehat{\widehat{\unsure{\mathcal C|\mathcal G}}}(\unsure{\Pi|\pi}'_0)$
is congruent modulo $M(\unsure{\Pi|\pi}'_0)$ to an element of
$\widehat{\widehat{\unsure{\mathcal C|\mathcal G}}}(\unsure{\Pi|\pi}_0)$
\margin{weaker than (b)}
--- and these conditions \unsure{together} imply that the canonical homomorphism
\diagram{\begin{tikzcd}[ampersand replacement=\&]
\mathbb{\Gamma}_{\deleted{\unsure{\pi}}\unsure{\Pi|\pi}_0}
  \arrow[r] \arrow[d, equal]
  \& \mathbb{\Gamma}_{\unsure{\Pi|\pi}'_0} \arrow[d, equal] \\
M(\unsure{\Pi|\pi}_0)/\widehat{\unsure{\mathcal C|\mathcal G}}(\unsure{\Pi|\pi}_0)
  \& M^{\unsure{\sharp|*}}(\unsure{\Pi|\pi}'_0)/\widehat{\unsure{\mathcal C|\mathcal G}}(\unsure{\Pi|\pi}'_0)
\end{tikzcd}}
is an isomorphism -- it even suffices for this to have
\unsure{(c)|(e)} (to be able to define this map) and (a)
(for its injectivity), and (b) and the strengthening (d) of
\unsure{(c)|(e)} (for its \unsure{bijectivity|surjectivity}. So we see that
\unsure{everything} hangs on the properties of the homomorphism of group inclusions
\diagram{\begin{tikzcd}[ampersand replacement=\&,column sep=2.2em,row sep=3.8em,
execute at end picture={
  \draw[line width=.35pt] ([xshift=-6pt,yshift=-1pt]\tikzcdmatrixname-1-3.south) -- ([xshift=-9pt,yshift=-1pt]\tikzcdmatrixname-1-3.south) -- ([xshift=-9pt,yshift=1pt]\tikzcdmatrixname-2-3.north) -- ([xshift=-6pt,yshift=1pt]\tikzcdmatrixname-2-3.north);
  \draw[line width=.35pt] ([xshift=17pt,yshift=-1pt]\tikzcdmatrixname-1-3.south) -- ([xshift=20pt,yshift=-1pt]\tikzcdmatrixname-1-3.south) -- ([xshift=20pt,yshift=1pt]\tikzcdmatrixname-2-3.north) -- ([xshift=17pt,yshift=1pt]\tikzcdmatrixname-2-3.north);
}]
{} \arrow[d,phantom,"{(7)}"{description}] \&
\widehat{\mathcal G}^{+}(\pi_0)
\arrow[r,phantom,"{\subset}"{description}]
\arrow[d,hook,no head] \&
M(\pi_0)
\arrow[r,phantom,"{\subset}"{description}]
\arrow[d,hook]
\arrow[d,phantom,"{?}"{description,xshift=8pt}] \&
\widehat{\widehat{\mathcal G}}(\pi_0) \arrow[d,hook,no head] \\
{} \&
\widehat{\mathcal G}^{+}(\pi'_0)
\arrow[r,phantom,"{\subset}"{description}] \&
M(\pi'_0)
\arrow[r,phantom,"{\subset}"{description}] \&
\widehat{\widehat{\mathcal G}}(\pi'_0)
\end{tikzcd}}
with the task of proving a) b) (which \unsure{concern} only the composite square)
and \unsure{c)|e)}.

\deleted{Besides d) \unsure{will be} an
\inserted{\unsure{immediate}} consequence of a), and of the fact
(which will ensure \unsure{c)|e)}) that every element of
$\widehat{\unsure{\mathcal C|\mathcal G}}^{+}(\unsure{\Pi|\pi}_0)$ is congruent modulo
$\unsure{\mathcal C|\mathcal G}$}

\scanpage{360}
\mspage{182}

Perhaps the properties a) b) c) are
completely false -- yet a) (which corresponds to the injectivity of
$P(\unsure{\Pi|\pi})\longrightarrow P(\unsure{\Pi|\pi}')$) seems fairly plausible.
Property b) (which would express the surjectivity of
$P(\unsure{\Pi|\pi})\longrightarrow P(\unsure{\Pi|\pi}')$) is much more problematic,
even doubtful. The trouble is that the group
$\widehat{\widehat{\unsure{\mathcal C|\mathcal G}}}(\unsure{\Pi|\pi}_0)$
can be ``much smaller'' than
$\widehat{\widehat{\unsure{\mathcal C|\mathcal G}}}(\unsure{\Pi|\pi}'_0)$,
we realize this by passing to the quotient by the subgroup
$\widehat{\unsure{\Pi|\pi}}'_0=\unsure{\Pi|\pi}'$ (contained in the smallest of the groups in (7),
and normal in the largest), we find on the second line
the groups $\widehat{\mathcal T}^{+}(\unsure{\Pi|\pi}'_0)$,
\deleted{\illegible{1}} $\unsure{\mathcal N|N}(\unsure{\Pi|\pi}'_0)$ and
$\widehat{\widehat{\mathcal T}}(\unsure{\Pi|\pi}'_0)$, on the first
extensions of the \unsure{corresponding|corresponding} groups
for $\unsure{\Pi|\pi}_0$ ($\widehat{\mathcal T}^{+}(\unsure{\Pi|\pi}_0)$,
$\unsure{\mathcal N|N}(\unsure{\Pi|\pi}_0)$ and
$\widehat{\widehat{\mathcal T}}(\unsure{\Pi|\pi}_0)$) by
\inserted{the finite group} $G=\unsure{\Pi|\pi}_0/\unsure{\Pi|\pi}'_0\simeq\unsure{\Pi|\pi}/\unsure{\Pi|\pi}'$,
\deleted{\unsure{let them be|let it be}
$\unsure{\widetilde{\widehat{\mathcal T}}^{+}|\widetilde{\mathcal T}^{+}}$
or $\mathcal T$}
\unsure{denoted|denoted} by $\sim$'s, so that \unsure{.|--}
\diagram{\begin{tikzcd}[ampersand replacement=\&,column sep=1.9em,row sep=3.8em,
execute at end picture={
  \draw[line width=.35pt] ([xshift=-6pt,yshift=-1pt]\tikzcdmatrixname-1-4.south) -- ([xshift=-9pt,yshift=-1pt]\tikzcdmatrixname-1-4.south) -- ([xshift=-9pt,yshift=1pt]\tikzcdmatrixname-2-4.north) -- ([xshift=-6pt,yshift=1pt]\tikzcdmatrixname-2-4.north);
  \draw[line width=.35pt] ([xshift=17pt,yshift=-1pt]\tikzcdmatrixname-1-4.south) -- ([xshift=20pt,yshift=-1pt]\tikzcdmatrixname-1-4.south) -- ([xshift=20pt,yshift=1pt]\tikzcdmatrixname-2-4.north) -- ([xshift=17pt,yshift=1pt]\tikzcdmatrixname-2-4.north);
}]
{} \arrow[d,phantom,"{(8)}"{description}] \&
G \arrow[r,phantom,"{\subset}"{description}] \arrow[dr,hook] \&
\widetilde{\widehat{\unsure{\mathcal T}}}^{+\unsure{:}}(\pi_0)
\arrow[r,phantom,"{\subset}"{description}]
\arrow[d,hook,no head] \&
\widetilde{\mathcal N}(\pi_0)
\arrow[r,phantom,"{\subset}"{description}]
\arrow[d,hook]
\arrow[d,phantom,"{?}"{description,xshift=8pt}] \&
\widetilde{\widehat{\widehat{\mathcal T}}}(\pi_0) \arrow[d,hook] \\
{} \& {} \&
\widehat{\mathcal T}^{+}(\pi'_0)
\arrow[r,phantom,"{\subset}"{description}] \&
\mathcal N(\pi'_0)
\arrow[r,phantom,"{\subset}"{description}] \&
\widehat{\widehat{\mathcal T}}(\pi'_0)
\end{tikzcd}},
which shows that the groups on the first line